\section{Implementação Concorrente --- Dados em Memória}

\subsection{Implementação em Java}

\subsubsection{Descrição}

A v3 parte de v2-virtual, a variante com melhor desempenho na seção anterior, e
elimina o principal gargalo identificado no profile: o parsing CSV repetido a cada
iteração. A mudança central é carregar todos os pontos do dataset em memória uma
única vez antes do loop de iterações.

\texttt{CsvDataset.loadAllPoints()} lê o arquivo e retorna um \texttt{double[][]}
com todos os pontos. O método \texttt{KMeans.cluster()} ganhou uma segunda assinatura
que recebe esse array diretamente, sem tocar em disco durante o clustering.
\texttt{KMeansBenchmark} pré-carrega o array no \texttt{@Setup(Level.Trial)},
excluindo a leitura do CSV do tempo medido. \texttt{KMeansSampler} usa um campo
\texttt{static volatile double[][] sharedPoints} com double-checked locking para que
todos os threads JMeter compartilhem uma única cópia dos dados em memória.

Em vez de criar lotes de \texttt{BATCH\_SIZE} pontos e coletar uma lista de
\texttt{Future}s, a v3 divide o array em $N$ chunks de tamanho fixo (onde $N$ é o
número de processadores lógicos) e usa \texttt{CountDownLatch(N)} para coordenar a
conclusão. Cada worker recebe índices \texttt{from}/\texttt{to} e opera diretamente
sobre o array compartilhado (somente leitura) com acumulação em arrays locais.

\begin{table}[H]
\centering
\begin{tabular}{ll}
\hline
\textbf{Componente} & \textbf{Alteração} \\
\hline
\texttt{CsvDataset}    & adicionado \texttt{loadAllPoints()} \\
\texttt{KMeans}        & nova assinatura \texttt{cluster(double[][], ...)};
                         \texttt{CountDownLatch} em vez de lista de \texttt{Future}s;
                         \texttt{processChunk} em vez de \texttt{processAssignBatch};
                         \texttt{DoubleAdder} para SSE \\
\texttt{KMeansBenchmark} & \texttt{double[][] points} pré-carregado em \texttt{@Setup(Level.Trial)} \\
\texttt{KMeansSampler}   & \texttt{static volatile} \texttt{sharedPoints} com double-checked locking \\
\hline
\end{tabular}
\caption{Alterações da v3 em relação à v2-virtual.}
\label{tab:v3-changes}
\end{table}

\subsubsection{Resultados}

\paragraph{Avaliação de Condição de Corrida}

Os mesmos quatro testes JCStress (JCS1--JCS4) da seção anterior foram executados
sobre a v3. A estrutura de acumulação local por worker foi mantida: cada chunk opera
em \texttt{localSums} e \texttt{localCounts} privados, sem compartilhamento.

\begin{table}[H]
\centering
\begin{tabular}{llll}
\hline
\textbf{ID} & \textbf{Resultado} & \textbf{Outcome FORBIDDEN} & \textbf{Freq.\ máx.} \\
\hline
JCS1 & \textbf{PASSOU}   & nenhum                                        & ---     \\
JCS2 & \textbf{PASSOU}   & nenhum                                        & ---     \\
JCS3 & FALHOU (esperado) & \texttt{"1"} lost update em \texttt{counts[0]++}   & 4{,}05\% \\
JCS4 & FALHOU (esperado) & \texttt{"1.0"} lost update em \texttt{sums[0]+=1.0} & 5{,}95\% \\
\hline
\end{tabular}
\caption{Resultados JCStress --- regiões críticas da v3.}
\label{tab:jcstress-v3}
\end{table}

JCS1 e JCS2 passam: leitura concorrente do array de centroides é segura. JCS3 e
JCS4 falham como esperado, confirmando que a race condition existe quando os arrays
são compartilhados --- o que a v3 evita da mesma forma que a v2.

\paragraph{Avaliação com Microbenchmark}

Os microbenchmarks foram executados com JMH 1.37 sobre JDK 25.0.1
(OpenJDK 64-Bit Server VM), com 1 fork, 1 iteração de aquecimento de 10 s e 3
iterações de medição de 10 s cada. O dataset foi pré-carregado em memória antes do
início das medições (\texttt{@Setup(Level.Trial)}).

\begin{table}[H]
\centering
\begin{tabular}{llrrrr}
\hline
\textbf{ID} & \textbf{Benchmark} & \textbf{Serial} & \textbf{v3} & \textbf{Erro ($\pm$)} & \textbf{Unidade} \\
\hline
B1 & Execução completa (\textit{k}=10, maxIter=20) & 212{,}57 &  4{,}55 & 17{,}98 & s/op \\
B2 & Custo por iteração (\textit{k}=10, maxIter=5) &  60{,}05 &  1{,}60 &  0{,}17 & s/op \\
B3 & \texttt{closestCentroid} (\textit{k}=10, 38 dims)  & 96{,}82 & 96{,}82 & --- & ns/op \\
B4 & \texttt{squaredDistance} (38 dims)                  &  9{,}31 &  9{,}31 & --- & ns/op \\
\hline
\end{tabular}
\caption{Resultados JMH --- serial vs.\ v3 (dados em memória).}
\label{tab:jmh-v3}
\end{table}

B2 cai de 60 s/op para 1{,}60 s/op: a eliminação do parsing CSV por iteração reduz
o custo em $\approx$37$\times$. B1 tem variância elevada ($\pm$18 s) por depender
do número de iterações até convergência. B3 e B4 são inalterados (computação pura).

\paragraph{Avaliação com Macrobenchmark}

Os macrobenchmarks foram executados com JMeter 5.6.3 sobre JDK 25.0.1, com
$k = 10$, \texttt{maxIter} = 3 e \texttt{tol} = 0{,}0. O dataset é carregado em
memória uma única vez por processo JMeter. Cada configuração rodou por 120 s com
ramp-up de 30 s, utilizando ParallelGC.

\begin{table}[H]
\centering
\begin{tabular}{rrrr}
\hline
\textbf{Threads JMeter} & \textbf{Execuções} & \textbf{Throughput (exec/min)} & \textbf{Latência média (ms)} \\
\hline
 1 & 105 & 52{,}5 &  1{.}145 \\
 2 & 112 & 56{,}0 &  2{.}010 \\
 4 & 116 & 57{,}1 &  3{.}796 \\
 8 & 119 & 57{,}2 &  7{.}457 \\
16 & 134 & 62{,}3 & 13{.}374 \\
\hline
\end{tabular}
\caption{Throughput e latência por número de threads JMeter --- v3, ParallelGC.}
\label{tab:jmeter-v3}
\end{table}

Com 1 thread JMeter, o throughput é 52{,}5 exec/min. O ganho com mais threads é
moderado (52,5 $\to$ 62,3 exec/min de 1 para 16 threads) porque cada execução já
satura todos os núcleos internamente. Nenhum erro foi registrado em nenhuma
configuração.

\paragraph{Avaliação de Profile}

O profile foi capturado com JFR (\texttt{settings=profile}) sobre uma execução
completa do algoritmo ($k=10$, \texttt{maxIter=3}, dataset $\approx$1{,}75 GB em
memória). A gravação foi analisada com a ferramenta \texttt{jfr} do JDK 25.

\textbf{I/O.}
Nenhum evento \texttt{jdk.FileRead} foi registrado durante o clustering. O arquivo
CSV é lido apenas na fase de inicialização; as iterações operam exclusivamente sobre
o array em memória.

\textbf{GC.}
Foram realizadas 60 coletas (51 Young G1New + 9 Old G1Old) com pausa total de
1.024 ms e pausa média de 17 ms por evento. O heap atingiu pico de 2{,}7 GB,
correspondente ao dataset carregado ($\approx$1{,}75 GB) mais overhead de trabalho.

\textbf{Threads.}
O pool de virtual threads usa os $N$ carriers do \texttt{ForkJoinPool} interno.
Chunks terminam em $\approx$100--200 ms, sendo facilmente visíveis ao JFR ao
contrário dos lotes de 5.000 pontos da v2.

\textbf{Síntese.}
O parsing CSV desaparece completamente do perfil de execução. O tempo é gasto
inteiramente em computação aritmética (\texttt{closestCentroid},
\texttt{squaredDistance}) distribuída entre os $N$ workers.

\iffalse % ============================================================
% ERLANG — seção desativada, manter para uso futuro
\subsection{\sout{Implementação em Erlang}}

\subsubsection{\sout{Descrição}}

\subsubsection{\sout{Resultados}}

\paragraph{\sout{Avaliação de Condição de Corrida}}

\paragraph{\sout{Avaliação com Microbenchmark}}

\paragraph{\sout{Avaliação com Macrobenchmark}}

\paragraph{\sout{Avaliação de Profile}}
\fi % ============================================================

\subsection{Discussão}

\textbf{Eliminação do gargalo de I/O.}
Nas versões v2, o parsing CSV consumia 86--99\% do CPU em cada iteração. Na v3,
nenhum evento de leitura de arquivo é registrado durante o clustering. O custo por
iteração (B2) caiu de $\approx$60 s para $\approx$1{,}6 s --- uma redução de 37$\times$.

\textbf{Comparação com versões anteriores --- microbenchmark.}

\begin{table}[H]
\centering
\begin{tabular}{lrrrr}
\hline
\textbf{Abordagem} & \textbf{B1 (s/op)} & \textbf{B2 (s/op)} & \textbf{Erro B1} & \textbf{Erro B2} \\
\hline
Serial      & 212{,}57 & 60{,}05 &  $\pm$34{,}48 &  $\pm$2{,}46 \\
v2-virtual  & 211{,}64 & 71{,}61 & $\pm$260{,}68 & $\pm$227{,}05 \\
v3          &   4{,}55 &  1{,}60 &  $\pm$17{,}98 &   $\pm$0{,}17 \\
\hline
\end{tabular}
\caption{Comparativo JMH: serial, v2-virtual e v3 (s/op).}
\label{tab:jmh-comparativo-v3}
\end{table}

A v3 é a primeira abordagem a superar o serial: B2 de 1{,}60 s contra 60{,}05 s,
uma melhora de $\approx$37$\times$. A variância de B1 ($\pm$18 s) é muito menor
do que em v2-virtual ($\pm$260 s), porque a carga de trabalho de cada iteração
é agora uniforme e previsível.

\textbf{Comparação com versões anteriores --- macrobenchmark.}

\begin{table}[H]
\centering
\begin{tabular}{lrrrrr}
\hline
\textbf{Abordagem} & \textbf{1 t} & \textbf{2 t} & \textbf{4 t} & \textbf{8 t} & \textbf{16 t} \\
\hline
v2-virtual (ParallelGC) & 1{,}62 & 2{,}13 & 2{,}14 & 3{,}77 &  4{,}34 \\
v3 (ParallelGC)         & 52{,}5 & 56{,}0 & 57{,}1 & 57{,}2 & 62{,}3  \\
\hline
\end{tabular}
\caption{Throughput (exec/min, ParallelGC) --- v2-virtual vs.\ v3.}
\label{tab:jmeter-comparativo-v3}
\end{table}

Com 1 thread JMeter, o throughput saltou de 1{,}62 para 52{,}5 exec/min
($\approx$32$\times$). Com 16 threads, de 4{,}34 para 62{,}3 exec/min
($\approx$14$\times$). O ganho menor em 16 threads ocorre porque múltiplos
processos de clustering competem pelos mesmos núcleos.

\textbf{Custo de memória.}
O heap pico aumentou de $\approx$250 MB (v2) para $\approx$2{,}7 GB (v3), reflexo
do dataset inteiro em memória. O número de eventos GC caiu (60 vs.\ 1.000+ na v2)
porque há muito menos alocação transiente por iteração.
